% ED26-METHOD-imprecise-probability mathematical component draft.
% Exact Walley2014 article acquisition remains open in the source note.
\chapter{Imprecise Probability and Lower-Prevision Reasoning}
\label{sec:p15-imprecise-probability}

A lower probability records a bound on an event; a lower prevision records
a bound on the value of a gamble. Coherence constrains their joint use.
It neither selects a uniquely true distribution nor requires every pair of
decisions to be ranked. The bounded-gamble construction here follows the
identified author manuscript of Miranda
\cite[author manuscript, §§2--4]{MirandaCoherentLowerPrevisionsCapture}.
The separation and finitely additive integration premises are proved in
\cref{thm:p17-mmeu-hahn-banach,thm:p17-mmeu-charges-compactness}.
The ambient set theory includes Choice as in that treatment.

\section{Gambles, prices and domains}
\label{sec:p15-ip-domains}

\begin{hbdefinition}{def:p15-ip-domain}{The bounded-gamble domain}
Let \(\Omega\ne\varnothing\). A gamble is a bounded real function on
\(\Omega\), with rewards expressed in a specified linear utility scale.
Write \(\mathcal L(\Omega)\) for all such functions.
More generally let \(\Sigma\) be an algebra of subsets of \(\Omega\),
and let \(B(\Sigma)\) be the uniform closure of its finite-valued simple
functions. Put \(B=B(\Sigma)\). It contains constants and event indicators
\(I_A\), \(A\in\Sigma\), and is closed under products.
For \(\Sigma=\mathcal P(\Omega)\) it is exactly \(\mathcal L(\Omega)\).
On a sigma algebra it is exactly the bounded measurable functions.
An assessment domain \(\mathcal K\subseteq B\) need not be linear,
finite or closed under negation.
\end{hbdefinition}

\begin{hbdefinition}{def:p15-ip-prices}{Lower and upper previsions}
A lower assessment \(\underline P:\mathcal K\to\mathbb R\) specifies
supremum acceptable buying prices. It licenses the transaction
\(f-\underline P(f)+\epsilon\) for each \(\epsilon>0\), without
specifying acceptance at the endpoint \(\underline P(f)\).
An upper assessment is an infimum selling price. The conjugate upper
prevision on \(-\mathcal K\) is
\[
 \overline P(-f)=-\underline P(f)\qquad(f\in\mathcal K).
\]
When all relevant indicators are in the domain, use
\(\underline P(A)=\underline P(I_A)\) and
\(\overline P(A)=\overline P(I_A)\).
Imprecision of a utility-valued price and imprecision of an event
probability are distinct inputs. For monetary rewards a linear utility
scale is an additional modelling assumption.
\end{hbdefinition}

\begin{hbdefinition}{def:p15-ip-coherence}{Avoiding sure loss and coherence}
Put \(G(f)=f-\underline P(f)\). The assessment avoids sure loss when
\[
 \sup_{\omega\in\Omega}\sum_{i=1}^n s_iG(f_i)(\omega)\geq0
 \quad(n\geq0,\ f_i\in\mathcal K,\ s_i\geq0).
\]
It is coherent when it avoids sure loss and also, for every such finite
portfolio, \(f_0\in\mathcal K\) and \(s_0\geq0\),
\[
 \sup_{\omega\in\Omega}
 \left[\sum_{i=1}^n s_iG(f_i)-s_0G(f_0)\right](\omega)\geq0.
\]
Only one assessed gamble is sold, possibly at multiple stakes. This
condition prevents the other accepted buying transactions from forcing
an increase of its stated supremum buying price. A linear prevision
instead permits arbitrary finite buying and selling portfolios.
\end{hbdefinition}

\begin{hblemma}{thm:p15-ip-stakes}{proved=annex:p15.ip.stakes}{Integer and real stakes}
The preceding conditions are unchanged if stakes are initially restricted
to nonnegative integers with repeated occurrences of a gamble.
\end{hblemma}

\section{Natural extension and compatible precise models}
\label{sec:p15-ip-extension}

\begin{hbdefinition}{def:p15-ip-natural-extension}{Natural extension of an assessment}
For \(g\in B\), define
\[
 \underline E(g)=\sup\left\{\alpha\in\mathbb R:
 g-\alpha\geq\sum_{i=1}^n s_i
             [f_i-\underline P(f_i)],\quad
 n\geq0,\ f_i\in\mathcal K,\ s_i\geq0\right\}.
\]
The inequality is pointwise and the empty portfolio is allowed.
Equivalently take the supremum over portfolios of
\(\inf_\Omega[g-\sum_i s_iG(f_i)]\).
\end{hbdefinition}

\begin{hbtheorem}{thm:p15-ip-natural-extension}{proved=annex:p15.ip.natural-extension}{Existence, coherence and least commitment}
If \(\underline P\) avoids sure loss, \(\underline E\) is finite,
coherent on \(B\), and the least coherent lower prevision on \(B\)
dominating \(\underline P\) on \(\mathcal K\). It satisfies
\[
\begin{gathered}
 \inf g\leq\underline E(g)\leq\sup g,\qquad
 \underline E(c)=c,\qquad
 \underline E(g+c)=\underline E(g)+c,\\
 \underline E(tg)=t\underline E(g)\ (t\geq0),\qquad
 \underline E(g+h)\geq\underline E(g)+\underline E(h),\\
 g\leq h\Longrightarrow\underline E(g)\leq\underline E(h),\qquad
 |\underline E(g)-\underline E(h)|\leq\|g-h\|_\infty.
\end{gathered}
\]
It agrees with the input on \(\mathcal K\) exactly when the input is
coherent. If the input incurs sure loss, \(\underline E(g)=+\infty\)
for every \(g\); it then supplies no finite coherent extension.
\end{hbtheorem}

\begin{hbtheorem}{thm:p15-ip-linear-domain}{proved=annex:p15.ip.linear-domain}{Coherence on the full vector domain}
A real functional \(I:B\to\mathbb R\) is a coherent lower prevision
if and only if
\[
 I(g)\geq\inf g,\qquad I(tg)=tI(g)\ (t>0),\qquad
 I(g+h)\geq I(g)+I(h).
\]
These conditions imply the displayed bounds, cash translation, monotonicity
and Lipschitz properties in \cref{thm:p15-ip-natural-extension}.
Oddness and additivity are additional
conditions; they do not follow from positive homogeneity.
\end{hbtheorem}

\begin{hbdefinition}{def:p15-ip-compatible}{Compatible linear previsions and probabilities}
A precise model is a positive normalized linear functional \(L\) on \(B\).
It corresponds to a finitely additive probability on \(\Sigma\), with
\(L(I_A)=L(A)\); its expectation is the uniform simple-function integral.
The compatible set is
\[
 \mathcal M(\underline P)=
 \{L:\ L\text{ is positive normalized linear on }B,
                   \ L(f)\geq\underline P(f)\ (f\in\mathcal K)\}.
\]
Use the eventwise, equivalently the bounded-integral, weak-star topology.
No countable additivity assumption has been inserted.
\end{hbdefinition}

\begin{hbtheorem}{thm:p15-ip-representation}{proved=annex:p15.ip.representation}{Representation on arbitrary assessment domains}
Avoidance of sure loss is equivalent to
\(\mathcal M(\underline P)\ne\varnothing\). In that case this set is
compact and convex, \(\mathcal M(\underline P)=\mathcal M(\underline E)\),
and
\[
 \underline E(g)=\min_{L\in\mathcal M(\underline P)}L(g)\qquad(g\in B).
\]
The input is coherent exactly when its values on \(\mathcal K\) equal
these minima. For a coherent functional on all of \(B\), its compatible
set is the unique nonempty closed convex representing set. This uniqueness
does not apply to an arbitrary proposed prior subset when only a restricted
assessment domain has been elicited.
\end{hbtheorem}

\begin{hbcorollary}{thm:p15-ip-dual}{proved=annex:p15.ip.dual}{Upper envelopes and event conjugacy}
Set \(\overline E(g)=-\underline E(-g)\). Then
\[
 \overline E(g)=\max_{L\in\mathcal M(\underline P)}L(g),\qquad
 \underline E(g)\leq\overline E(g),\qquad
 \overline E(A)=1-\underline E(A^c).
\]
Natural extension through an intermediate domain has the same final
answer as direct extension, provided that intermediate assessments are
the natural-extension values.
\end{hbcorollary}

For the remainder of the chapter assume that the input avoids sure loss,
so that \(\underline E\) and its nonempty compatible set are defined
with the preceding properties.

\section{Finite additivity and countable additivity}
\label{sec:p15-ip-additivity}

\begin{hbtheorem}{thm:p15-ip-countable-additivity}{proved=annex:p15.ip.countable-additivity}{The extra continuity condition}
Suppose \(\Sigma\) is a sigma algebra. Every member of
\(\mathcal M(\underline P)\) is countably additive if and only if,
for every sequence \(A_n\downarrow\varnothing\),
\[
                         \overline E(A_n)\longrightarrow0.
\]
Coherence alone does not imply this condition. On a countably infinite
space a free-ultrafilter probability is coherent but not countably
additive. Even \(\underline E(A_n)\to0\) for every such sequence does
not suffice: the vacuous lower prevision \(\underline E(g)=\inf g\)
has that lower continuity and admits non-countably-additive precise models.
\end{hbtheorem}

Requiring countably additive compatible models is therefore a substantive
extra assumption. It can change compactness, attainment and the available
models. The indicator of a conditioning event must belong to the selected
event algebra; an undefined probability of a set outside that algebra is
not a zero probability.

\section{Event bounds do not determine all gamble bounds}
\label{sec:p15-ip-event-information}

\begin{hbtheorem}{thm:p15-ip-event-separation}{proved=annex:p15.ip.event-separation}{Equal event envelopes, different coherent gamble prices}
On \(\Omega=\{1,2,3\}\), let \(C\) be the convex hull of all six
permutations of \((1,3,6)/10\), and let \(D\) be the hull with
\((1,3,6)/10\) omitted. Their lower probabilities agree on every event:
zero on the empty event, \(1/10\) on singletons, \(2/5\) on pairs,
and one on the whole space. Their upper probabilities also agree.
But for \(f=(2,1,0)\),
\[
                \min_C E_p f=1/2,\qquad \min_D E_p f=7/10.
\]
The decision against the certain payoff \(3/5\) consequently changes
under a stated maxmin rule. The missing gamble information cannot be
recovered by simply assigning a precise number within each event interval.
\end{hbtheorem}

Belief/plausibility and possibility/necessity are separately defined event
constructions in \cref{def:p14-evidence-uncertainty-data}. Under its finite
mass-allocation assumptions, \cref{thm:p14-evidence-finite-measures}
supplies the belief/plausibility allocated probability envelopes. Taking expected
values under those probabilities also supplies coherent lower previsions
by \cref{thm:p15-ip-representation}; the special model's assumptions
determine which probabilities are retained. The explicit coherent lower
probability which fails the belief-function inequality in
\cref{thm:p14-evidence-measure-separation} proves that general coherent
lower probabilities need not be belief functions. A possibility union law
uses a maximum, rather than the additive law of a probability.
These representations cannot be identified solely by the common range
\([0,1]\), and their conditioning or evidence-combination rules require
their own hypotheses.
\Cref{sec:p15-w02b-p15-034-evidence-models} proves that the gamble
extension of a belief function is the focal-set minimum sum, equal to its
Choquet value (\cref{thm:p15-w02b-p15-034-belief-extension}); that a
coherent lower prevision can agree with a belief function on every event
and still differ from that extension
(\cref{thm:p15-w02b-p15-034-belief-not-prevision}); that necessity
measures form a proper subclass of belief functions
(\cref{thm:p15-w02b-p15-034-strict-classes}); and that Dempster
conditioning of a belief function or a possibility measure need not be
jointly coherent with the natural extension
(\cref{thm:p15-w02b-p15-034-dempster}).

\begin{hblemma}{thm:p15-ip-possibility-bridge}{proved=annex:p15.ip.possibility-bridge}{A normalized finite possibility model as a coherent envelope}
For the finite normalized possibility distribution in
\cref{def:p14-evidence-uncertainty-data}, order the outcomes so that
\(1=\pi_1\geq\cdots\geq\pi_n\geq0\), and put
\(\pi_{n+1}=0\). The nested focal sets
\(B_i=\{\omega_1,\ldots,\omega_i\}\), with masses
\(m(B_i)=\pi_i-\pi_{i+1}\), have
\[
 \operatorname{Bel}_m(A)=\operatorname{Nec}_\pi(A),\qquad
 \operatorname{Pl}_m(A)=\operatorname{Poss}_\pi(A).
\]
Their allocated probabilities therefore give coherent lower/upper gamble
previsions extending these particular event bounds. The nested-focal-set
constraint is an additional restriction on the general model.
\end{hblemma}

\section{Conditional models for one observation event}
\label{sec:p15-ip-conditioning}

\begin{hbdefinition}{def:p15-ip-conditional}{Conditional prices and compatible pairs}
Let \(A\in\Sigma\), \(A\ne\varnothing\). A precise conditional
model \(Q\) is a positive normalized linear functional on \(B\),
supported on \(A\): \(Q(I_A)=1\). A pair \((L,Q)\) is Bayes-compatible
when
\[
                   L(I_Ag)=L(A)Q(g)\qquad(g\in B).
\]
For \(L(A)>0\) this determines \(Q\); for \(L(A)=0\) it constrains
none of its supported values. A conditional transaction has payoff
\(I_A[g-q]\), so it is called off outside \(A\).
Coherence of transactions made only on \(A\) is assessed on \(A\),
rather than being made automatic by the zero payoff outside it.
Joint coherence here tests every finite portfolio of unconditional and
conditional buying gains with at most one assessed gain sold, using
nonnegative stakes. The supremum is taken on the union of their active
conditioning events: on \(\Omega\) if an unconditional transaction is
present, and otherwise on \(A\). The separate conditional tests on
\(A\) are required as well.
These definitions concern one event. They do not assert a simultaneous
updating theorem for arbitrary infinite partitions.
\end{hbdefinition}

\begin{hbdefinition}{def:p15-ip-updates}{Natural and regular conditioning}
The natural conditional lower prevision is the infimum of \(Q(g)\)
over all Bayes-compatible pairs with \(L\in\mathcal M(\underline P)\).
For \(\overline E(A)>0\), the regular extension is
\[
 \underline R(g\mid A)=
 \inf_{\substack{L\in\mathcal M(\underline P)\\L(A)>0}}
                         \frac{L(I_Ag)}{L(A)}.
\]
If \(\overline E(A)=0\), define it to be \(\inf_A g\); this case
contains no positive-probability Bayesian conditional. Upper conditional
previsions are the conjugates on \(A\).
\end{hbdefinition}

Interpreting a contingent price as a future price after learning \(A\)
requires an updating principle relating transactions accepted after that
information to their called-off transactions before it, with no additional
information from the observation mechanism. This is a modelling premise,
rather than a consequence of the numerical bounds alone
\cite[author manuscript, §3.1]{MirandaCoherentLowerPrevisionsCapture}.
At lower probability zero, coherence permits several conditional models.
Choosing regular extension adds a commitment by excluding the zero-\(A\)
precise models from that conditional calculation; the displayed formulas
do not claim that coherence alone uniquely chooses this updating rule.

\begin{hbtheorem}{thm:p15-ip-conditional-models}{proved=annex:p15.ip.conditional-models}{The three probability boundaries}
Natural and regular conditional lower previsions are coherent on the
gamble space restricted to \(A\). If \(\underline E(A)>0\), they agree,
their Bayesian infimum is attained, and their value is determined by the
original assessment. If \(\underline E(A)=0<\overline E(A)\), natural
conditioning is vacuous, \(\inf_A g\), while regular conditioning can
be more informative. If \(\overline E(A)=0\), both are vacuous under
the stated convention. On this single-event domain the conditional models
are jointly coherent with the original unconditional assessment. Regular
conditioning, when \(\overline E(A)>0\), preserves its unconditional
envelope by the density of
positive-\(A\) models in the original compatible set.
\end{hbtheorem}

\begin{hbtheorem}{thm:p15-ip-generalized-bayes}{proved=annex:p15.ip.generalized-bayes}{Generalized Bayes rule and its zero boundary}
For \(\underline E(A)>0\), the equation
\[
                       \underline E(I_A[g-q])=0
\]
has exactly one root, namely \(q=\underline R(g\mid A)\).
For \(\underline E(A)=0<\overline E(A)\), its roots are every
\(q\leq\underline R(g\mid A)\); choosing an unspecified root then
does not define an update. Its largest root is regular extension.
\end{hbtheorem}

\begin{hbtheorem}{thm:p15-ip-conditional-bounds}{proved=annex:p15.ip.conditional-bounds}{What joint coherence actually forces at one event}
Let a full-domain conditional lower prevision \(C(\cdot\mid A)\)
be separately coherent and jointly coherent with \(\underline E\)
under \cref{def:p15-ip-conditional}. Then for every \(g\in B\),
\[
              \underline E(I_A[g-C(g\mid A)])=0.
\]
If \(\underline E(A)>0\), it follows that \(C=\underline R\).
If \(\underline E(A)=0<\overline E(A)\), its pointwise bounds are
\[
                    \inf_A g\leq C(g\mid A)\leq\underline R(g\mid A).
\]
Both limiting functionals are jointly coherent, so natural conditioning
is the least and regular conditioning the greatest on this single-event
domain. This conclusion does not select which intermediate coherent
assessments the subject actually supplies.
\end{hbtheorem}

\begin{hbexample}{def:p15-ip-zero-example}{A nonvacuous regular update at lower probability zero}
On \(\{a,b,c\}\) retain all distributions satisfying \(p_c\geq p_a\).
For \(A=\{a,c\}\), lower and upper probability are zero and one.
Natural conditioning gives lower probability zero to \(\{c\}\),
while regular conditioning gives \(1/2\): every retained positive-\(A\)
distribution has \(p_c/(p_a+p_c)\geq1/2\), with equality possible.
The prior concentrated at \(b\) explains the natural vacuity.
This is a calculation under specified models, not permission to divide
by the lower probability of \(A\).
It is the source's finite Example 1
\cite[author preprint, p.~2]{DeBockDeCooman2015ConditioningPreprint}.
\end{hbexample}

\section{Decision comparisons and independence assumptions}
\label{sec:p15-ip-decision-independence}

To assess robust improvement of \(f\) over \(g\), calculate
\(\underline E(f-g)\), as in \cref{def:p17-robust-dominance}.
Comparing \(\underline E(f)\) and \(\underline E(g)\) is the separate
maxmin rule. Complete utility information, a feasible menu and a stated
choice criterion are still needed; coherence alone does not specify them.
The vector example in \cref{thm:p15-ip-event-separation} uses common
known utility values, so its changed comparison is due to the retained
gamble information rather than a changed utility scale.

\begin{hbdefinition}{def:p15-ip-independence}{Specified notions of independence}
Epistemic irrelevance of \(Y\) for \(X\) requires that the conditional
lower prevision of every \(X\)-gamble after learning \(Y=y\) equals
its original marginal lower prevision. Epistemic independence requires
irrelevance in both directions. These are statements about the specified
conditional model. A strong product instead takes the closed convex hull
of all products of compatible marginal precise models, as in
\cref{def:p17-mmeu-product,thm:p17-mmeu-product-priors}.
The factorization concerns the generating products and their extreme-point
characterization, not every convex mixture. Logical independence of the
coordinates additionally means that the full product outcome space is
possible.
\end{hbdefinition}

\begin{hbtheorem}{thm:p15-ip-independence-boundary}{proved=annex:p15.ip.independence-boundary}{A product model does not settle null-event conditional prices}
Let \(X,Y\in\{0,1\}\). Retain the convex hull of product distributions
with \(p_X(1)\in[2/5,3/5]\), \(p_Y(1)\in[0,1]\).
The marginal lower probability of \(X=1\) is \(2/5\).
For \(A=\{Y=1\}\), its lower probability is zero, natural conditioning
gives \(\underline E(X=1\mid A)=0\), and regular conditioning gives
\(\underline R(X=1\mid A)=2/5\). Therefore any assertion of
irrelevance at the zero-lower-probability history must state which
conditional model, or which additional conditional assessments, it uses.
\end{hbtheorem}

A decision comparison that needs an independence judgement, with no
conditioning event of lower probability zero, is worked in
\cref{sec:p15-w02b-p15-034-independence-decision}: with imprecise binary
marginals, the agreement bet is robustly preferred to a sure payment under
the strong product, but the two acts stay unranked under the least
informative epistemically independent model and under the marginals alone
(\cref{thm:p15-w02b-p15-034-independence-models,thm:p15-w02b-p15-034-independence-decision}).

For infinite partitions, unconditional coherence does not by itself
justify simultaneous conditional assessments on every cell.
\Cref{sec:p15-w02b-p15-034-partition-updates} declares the
partition-update semantics, proves that a coherent partition update exists
exactly when the unconditional model is conglomerable for the partition
(\cref{thm:p15-w02b-p15-034-existence}), works a finitely additive
prevision with no partition update
(\cref{thm:p15-w02b-p15-034-counterexample}), and shows that the same
cellwise assessments pass Williams' finite-support criterion
(\cref{thm:p15-w02b-p15-034-williams}). The single-event results above
preserve their stated domain.
